\begin{abstract}
\noindent
Prediction markets aggregate dispersed beliefs into a single forecast, but their
incentives normally rely on a verifiable outcome: agents are eventually paid based
on what actually happened. Many decisions of interest --- the causal effect of a
policy, whether a piece of content violates a rule, the long-run consequence of a
choice --- never resolve cleanly, so this lever is unavailable. \textcite{srinivasan2025selfresolving}
propose a \emph{self-resolving} prediction market that removes the outcome
entirely: agents report sequentially, the market stops at a random time, and every
agent is paid the cross-entropy of their report against a late \emph{reference}
agent rather than against the truth. Their key insight is that a reference with
access to enough conditionally independent signals becomes a good proxy for the
ground truth, which makes truthful reporting a (strict) perfect Bayesian
equilibrium. Their guarantees are, however, entirely theoretical. This paper
reproduces the mechanism in an open-source agent-based simulator and stress-tests
it along two axes. First, in a clean synthetic world where the paper's assumptions
hold by construction, we estimate a focal agent's profit-maximising deviation by
Monte Carlo and confirm that the best response converges to truthful reporting as
the reference's number of informational substitutes $k$ grows ($\gamma^\star \to 1$
by $k \approx 64$), operationalising the paper's main theorems; we also reproduce
the analytic minimum-substitutes bound $k_{\min}$. Second, we run a counterfactual
study on roughly $90$ resolved Polymarket markets, treating each price path as a
stream of reports and paying it both ways --- self-resolving (against a late
reference price) and verifiable (against the realised outcome). The two payment
schemes agree almost perfectly when the reference is genuinely late (Pearson
$\approx 1.00$, payout ratio $\approx 0.95$) and degrade gracefully as the reference
is moved earlier, isolating the outcome-leakage channel that the reference relies
on. Together the results give empirical content to a purely theoretical mechanism
and map out where reference-based resolution succeeds and where it breaks.
\end{abstract}
